\documentclass[11pt]{article}

\usepackage[margin=25mm]{geometry}

\usepackage{amsmath,amssymb,fontspec,unicode-math}

\IfFontExistsTF{DejaVu Serif}{\setmainfont{DejaVu Serif}}{\setmainfont{Latin Modern Roman}}

\IfFontExistsTF{Latin Modern Math}{\setmathfont{Latin Modern Math}}{\setmathfont{DejaVu Math TeX Gyre}}

\usepackage{xurl,hyperref}

\setlength{\parindent}{0pt}\setlength{\parskip}{6pt}

\begin{document}

\title{Novel Exploration via Orthogonality}

\author{ELVIA preliminary AI analysis}\date{21 October 2026}\maketitle

Preliminary assessment: Major.\par

Independent scientific review remains pending. Author adoption has not been requested.

Strict eigenvector-magnitude ascent can terminate on an exact plateau before reaching the required novelty threshold. The undirected random-walk geometry also requires degree weighting, while a separate selector branch differs from the documented exploration protocol.

\section*{Review scope}

Claim-focused reading of the original paper and relevant author-released passages; direct mathematical or static-source reasoning. No author program, earlier verifier, model or benchmark was executed during this check.

\section*{Problem 1: Strict magnitude ascent can stop at a plateau}

Source: Theorem 3.3 and proof, PDF p.4; option.take\_option\_action

Assessment: supported; impact: major

\subsection*{Source claim (paraphrase)}

For \(L=I+\operatorname{diag}(N)-P\), starting at \(N_i>\operatorname{Re}\lambda\), the theorem promises a strictly increasing magnitude path ending below that threshold.

\subsection*{Analysis}

Use the directed cycle \(0\to1\to2\to0\), \(N=(1/4,1/8,11/28)\), \(f=(1,1,7/8)\), and \(\lambda=1/4\). Direct substitution gives \(Lf=\lambda f\), and the positive eigenvector is the principal mode after a nonnegative irreducible shift. Starting at \(2\), magnitude rises to \(1\) at node \(0\), but \(N_0=\lambda\). The only next step reaches node \(1\) at equal magnitude, violating strict ascent; only that plateau traversal reaches \(N_1<\lambda\). A connected undirected tree with edges \(01,12,23,14,45\), \(N=(1/2,1/2,0,0,0,0)\), \[f=(0,0,1/\sqrt2,1,-1/\sqrt2,-1)\], and \(\lambda=1-1/\sqrt2\) gives a second exact eigenpair: node \(0\) has only the zero-magnitude successor \(1\), so no first strict step exists.

\subsection*{Finding and scope}

Both positive plateaus and zero plateaus defeat the universal strict claim and the corresponding stop-on-ties interpretation. The more restricted source-maximum theorem and reported reward curves are not refuted.

\subsection*{Proposed correction}

With nonzero starting entry, strict ascent guarantees only an endpoint \(N\le\operatorname{Re}\lambda\). Allowing nondecreasing steps can restore the strict novelty threshold under strong connectivity, but requires plateau traversal with a visited set and changes the existing termination policy.

A proposed correction is retained where stated; independent verification and author adoption remain unestablished.

\section*{Problem 2: An undirected random walk is not Euclidean Hermitian}

Source: Section 3, PDF pp.3–4; Appendix D, PDF p.22

Assessment: supported; impact: major

\subsection*{Source claim (paraphrase)}

The operator \(I+\Gamma-D^{-1}A\) is called Hermitian on undirected graphs and assigned an unweighted Euclidean edge-energy formula.

\subsection*{Analysis}

For the three-vertex undirected path, degrees are \((1,2,1)\), so \(P_{01}=1\ne P_{10}=1/2\). With zero potential and \(x=(0,1,0)\), \(x^\top(I-P)x=1\), but the printed edge sum \(\tfrac12\sum_{ij}P_{ij}(x_i-x_j)^2=3/2\). Right eigenvectors \((1,1,1)\) and \((1,-1,1)\) have distinct eigenvalues \(0,2\) yet Euclidean inner product \(1\).

\subsection*{Finding and scope}

Euclidean symmetry and orthogonality fail on irregular graphs. Real eigenvalues survive by similarity to a symmetric matrix; regular graphs remain a valid special case.

\subsection*{Proposed correction}

Use the degree-weighted inner product and quotient \[\frac{\tfrac12\sum_{ij}A_{ij}(x_i-x_j)^2+\sum_i d_iN_ix_i^2}{\sum_i d_ix_i^2},\] or the symmetric operator \(D^{1/2}LD^{-1/2}\). This does not itself prove orthogonality of nonlinear magnitude-based policies.

A proposed correction is retained where stated; independent verification and author adoption remain unestablished.

\section*{Problem 3: Random-option selection occurs outside the epsilon event}

Source: Section 5, PDF p.7; all\_functions.py, AgentQ2.option\_action, lines 487–514

Assessment: supported; impact: minor

\subsection*{Source claim (paraphrase)}

The paper chooses a random option with probability initP conditional on an epsilon exploration step. The release makes a new option draw after the epsilon-greedy action.

\subsection*{Analysis}

With one eligible option, no active option, a greedy primitive action and independent draws \(u,v\), let \(\epsilon=\mathrm{initP}=0.1\). At \(u=0.5,v=0.05\), the release overrides a greedy action with an option; the described conditional procedure does not. The eligible-decision branch probability is \(0.1\), versus \(0.1\cdot0.1=0.01\). These are branch probabilities, not fractions of all environment transitions. An outer false exploration flag suppresses this direct branch, but recursive termination calls need not preserve that flag.

\subsection*{Finding and scope}

The released selector and described protocol differ. Published-run correspondence and reward-curve effects remain unestablished.

\subsection*{Proposed correction}

Either condition option sampling on the single epsilon event or document the independent override, according to the actual run protocol. No fully integrated selector repair is claimed.

A proposed correction is retained where stated; independent verification and author adoption remain unestablished.

\section*{Limits}

Preliminary AI assessment; independent human scientific adjudication remains pending.\par

Analytic witnesses and static release checks do not remeasure the published benchmark results.\par

\section*{Source identity}

NeurIPS 2025 official proceedings PDF (retained local copy)

\url{https://proceedings.neurips.cc/paper_files/paper/2025/file/bc5d44c37a9ccdf9300637a1d79ad9b8-Paper-Conference.pdf}

Reviewed PDF SHA-256: \nolinkurl{caf07f036239850d0a062a52b8319ebd781df0e68724c773c58f064a2954c899}

The SHA-256 identifies the exact retained PDF inspected in this rapid source/argument check. Original supplied report bytes are retained locally. No unambiguous original PDF retrieval timestamp was recovered; file modification time was not substituted.

\end{document}
